\begin{example}
Consider: 
$$
A = \begin{bmatrix}
    -1 & 1 \\
    -1 & 1
\end{bmatrix} \quad
\text{and} \quad
I_2 = \begin{bmatrix}
    1 & 0 \\
    0 & 1
\end{bmatrix}
$$
We wish to verify that: $|A+I_2| \leq |A|+|I_2|$ is not true. It is noted that $\sqrt{I_2}=I_2$. Using Theorem \ref{theorem1} allows simple computation of the matrix absolute value for matrix $A+I_2$ and matrix $A$.
\begin{align*}
    |A+I_2| &= 
    \left|\begin{bmatrix}
        -1 & 1 \\
        -1 & 1
    \end{bmatrix} 
    +
    \begin{bmatrix}
        1 & 0 \\
        0 & 1
    \end{bmatrix}\right| 
    = 
    \left|\begin{bmatrix}
        0 & 1 \\
        -1 & 2
    \end{bmatrix}\right| \\
    &=
    \left(
    \begin{bmatrix}
        0 & -1 \\
        1 & 2
    \end{bmatrix}
    \begin{bmatrix}
       0 & 1 \\
        -1 & 2 
    \end{bmatrix}\right)^\frac{1}{2} 
    =
    \begin{bmatrix}
        1 & -2 \\
        -2 & 5
    \end{bmatrix}^\frac{1}{2}\\
    &= 
    \frac{1}{\sqrt{8}}
    \begin{bmatrix}
        2 & -2\\
        -2 & 6
    \end{bmatrix}\\
    &=
    \begin{bmatrix}
         \frac{\sqrt{2}}{2} & \frac{-\sqrt{2}}{2} \\
        \frac{-\sqrt{2}}{2} & \frac{3\sqrt{2}}{2}
    \end{bmatrix}
\end{align*}
Similarly, 
\begin{align*}
    |A| &= \left|
    \begin{bmatrix}
        -1 & 1\\
        -1 & 1
    \end{bmatrix} \right|\\
    &=
    \left(
    \begin{bmatrix}
        -1 & -1\\
        1 & 1
    \end{bmatrix}
    \begin{bmatrix}
        -1 & 1\\
        -1 & 1
    \end{bmatrix}\right)^\frac{1}{2} 
    = 
    \begin{bmatrix}
        2 & -2 \\
        -2 & 2
    \end{bmatrix}^\frac{1}{2}\\
    &= 
    \frac{1}{2}
    \begin{bmatrix}
        2 & -2\\
        -2 & 2
    \end{bmatrix}\\
    &=
    \begin{bmatrix}
        1 & -1\\
        -1 & 1
    \end{bmatrix}
\end{align*}
If we apply the triangle inequality with these three matrices found, we get the following result:
\begin{align*}
    |A+I_2| &\leq |A| +I_2 \\
    \implies\begin{bmatrix}
        \frac{\sqrt{2}}{2} & \frac{-\sqrt{2}}{2} \\
        \frac{-\sqrt{2}}{2} & \frac{3\sqrt{2}}{2}
    \end{bmatrix}
    &\leq
    \begin{bmatrix}
        1 & -1 \\
        -1 & 1
    \end{bmatrix}
    +
    \begin{bmatrix}
        1 & 0 \\
        0 & 1
    \end{bmatrix}\\
    \Leftrightarrow 
    \begin{bmatrix}
        0 & 0 \\
        0 & 0
    \end{bmatrix}
    &\leq
    \begin{bmatrix}
        2 &-1\\
        -1 & 2
    \end{bmatrix} 
    -
    \begin{bmatrix}
        \frac{\sqrt{2}}{2} & \frac{-\sqrt{2}}{2} \\
        \frac{-\sqrt{2}}{2} & \frac{3\sqrt{2}}{2}
    \end{bmatrix}\\
    \implies
    \begin{bmatrix}
        0 & 0 \\
        0 & 0
    \end{bmatrix}
    &\leq
    \begin{bmatrix}
        2 - \frac{\sqrt{2}}{2} & -1 + \frac{\sqrt{2}}{2}\\
        -1 + \frac{\sqrt{2}}{2} & 2 - \frac{3\sqrt{2}}{2}
    \end{bmatrix} 
\end{align*}
If the difference matrix $|A|+I_2-|A+I_2|$ is not positive semidefinite then the triangle inequality would have failed.
We will check the determinants of the minor diagonals for this difference matrix to determine this.
\begin{align*}
    \left|2-\frac{\sqrt{2}}{2}\right| &\approx 2 - 0.707 \\
    & = 1.23 > 0
    \text{, and}\\
    \begin{vmatrix}
         2 - \frac{\sqrt{2}}{2} & -1 + \frac{\sqrt{2}}{2}\\
        -1 + \frac{\sqrt{2}}{2} & 2 - \frac{3\sqrt{2}}{2}
    \end{vmatrix}
    &= \left( 2 - \frac{\sqrt{2}}{2}\right)\left(2 - \frac{3\sqrt{2}}{2}\right) - \left(-1 + \frac{\sqrt{2}}{2}\right)^2\\
    &= \left(4 - 3\sqrt{2} - \sqrt{2} + \frac{3}{2} - \left(1 - \sqrt{2} + \frac{1}{2}\right)\right)\\
    &= 4 - 3\sqrt{2}\\
    &\approx 4 - 4.243\\
    &= -0.234 < 0.
\end{align*}
This shows that the difference matrix is not PSD; consequently, the triangle inequality fails given the choice of $A$ and $I_2$.
\end{example}